% !TEX root = ./article/chapter2-article.tex

This chapter develops the main technical apparatus of the
dissertation.  Building on the neighborhood semantics introduced in
Chapter~\ref{ch:chapter1}, we specialize the framework to deontic
logic and extend it to conditional operators.

\section{Deontic Neighborhood Frames}

We now interpret the modal operator as obligation.  In place of the
generic $\nec$, we write $\ought$ for ``it ought to be that'' and
$\may$ for ``it is permissible that,'' where $\may A \coloneqq
\neg\ought\neg A$.

\begin{definition}\label{def:deontic-nbhd-frame}
  A \textit{deontic neighborhood frame} is a neighborhood frame
  $\FN = \langle W, \fn \rangle$ where $\fn(w)$ is interpreted as
  the collection of obligation-neighborhoods at $w$.
\end{definition}

Frame conditions on $\fn$ correspond to axioms of the logic:

\begin{proposition}\label{prop:D-condition}
  A neighborhood frame validates $\dax$ (i.e., $\ought A \to \may
  A$) if and only if $\fn$ satisfies:
  \[
    \text{for all } w \in W,\ \text{if } X \in \fn(w) \text{ then }
    W \setminus X \notin \fn(w).
  \]
\end{proposition}

\begin{proof}
  $(\Rightarrow)$\quad Suppose the frame validates $\dax$ and $X
  \in \fn(w)$.  Let $V(p) = X$.  Then $w \models \ought p$, so by
  $\dax$, $w \models \may p$, i.e., $w \not\models \ought \neg p$.
  Hence $\truthset{\neg p} = W \setminus X \notin \fn(w)$.

  $(\Leftarrow)$\quad Assume the frame condition holds and $w
  \models \ought A$.  Then $\truthset{A} \in \fn(w)$, so $W
  \setminus \truthset{A} = \truthset{\neg A} \notin \fn(w)$, giving
  $w \not\models \ought \neg A$, i.e., $w \models \may A$.
\end{proof}


\section{Monotonicity and Congruentiality}

A key structural distinction is between \textit{monotone} and
merely \textit{congruential} logics.

\begin{definition}\label{def:monotone}
  A neighborhood function $\fn$ is \textit{supplemented}
  (monotone) if for all $w \in W$:
  \[
    X \in \fn(w) \text{ and } X \subseteq Y \implies Y \in \fn(w).
  \]
\end{definition}

\begin{note}\label{note:supplementation}
  Supplementation corresponds to the monotonicity rule (RM): from
  $\vdash A \to B$ infer $\vdash \ought A \to \ought B$.  Logics
  that validate the congruence rule \rerule\ --- from $\vdash A \iff
  B$ infer $\vdash \ought A \iff \ought B$ --- without validating RM
  are called \textit{congruential} or \textit{classical}.
\end{note}

\begin{lemma}\label{lem:supplementation}
  Every neighborhood function $\fn$ has a unique smallest monotone
  extension $\fn^+$, defined by
  \[
    \fn^+(w) = \{Y \subseteq W : X \subseteq Y \text{ for some }
    X \in \fn(w)\}.
  \]
\end{lemma}

\begin{proof}
  Clearly $\fn(w) \subseteq \fn^+(w)$.  To see that $\fn^+$ is
  supplemented, suppose $Y \in \fn^+(w)$ and $Y \subseteq Z$.  Then
  there exists $X \in \fn(w)$ with $X \subseteq Y \subseteq Z$, so
  $Z \in \fn^+(w)$.  Minimality follows by construction.
\end{proof}


\section{Conditional Operators}

For dyadic deontic logics, we move to a conditional obligation
operator $\cobs$, read ``given $A$, it ought to be that $B$.''

\begin{definition}\label{def:dyadic-nbhd-frame}
  A \textit{dyadic neighborhood frame} is a structure $\langle W,
  \fd, V \rangle$ where $\fd : W \times \wp(W) \to \wp(\wp(W))$ is
  a selection function mapping a world and a proposition to a
  collection of neighborhoods.
\end{definition}

The truth clause for conditional obligation is:
\[
  \MN, w \models \cobs(A, B) \iff \truthset{B} \in \fd(w,
  \truthset{A}).
\]

\begin{theorem}\label{thm:cond-correspondence}
  A dyadic neighborhood frame validates $\cobs(A, B \land C) \to
  \cobs(A, B)$ if and only if $\fd$ is monotone in the consequent.
\end{theorem}

\begin{proofsketch}
  The proof follows the same pattern as
  Proposition~\ref{prop:D-condition}: the left-to-right direction
  extracts the frame condition via a suitable valuation, and the
  right-to-left direction verifies the axiom directly.
\end{proofsketch}
